\documentclass[a4paper,oneside,11pt]{amsart}

%\usepackage[spanish, activeacute]{babel}
\usepackage[utf8]{inputenc}
%\usepackage[latin1]{inputenc}
\usepackage{latexsym}
\usepackage{multicol}
\usepackage{enumerate}
\usepackage{enumitem}
\usepackage{booktabs}
\usepackage[heightrounded]{geometry}
\geometry{top=2cm,bottom=2cm,left=2cm,right=2cm}
\usepackage{graphicx}

\DeclareMathOperator{\cond}{cond}
\DeclareMathOperator{\Id}{Id}
\DeclareMathOperator{\re}{Re}
%\DeclareMathOperator{\im}{Im}
\newcommand{\I}{\mathbb{I}}
\newcommand{\R}{\mathbb{R}}
\newcommand{\Q}{\mathbb{Q}}
\newcommand{\C}{\mathbb{C}}
\newcommand{\Z}{\mathbb{Z}}
\newcommand{\N}{\mathbb{N}}
\DeclareMathOperator{\im}{Im}

\newcommand{\nombremateria}
{Investigaci\'on Operativa}
\newcommand{\nombrecuatrimestre}
    {Segundo cuatrimestre de 2022}
\newcommand{\numeroparcial}
    {Recuperatorio del Segundo Parcial}
\newcommand{\fecha}
    {14 de Diciembre de 2022}
\newcommand{\numerotema}
    {\bf 1}

%\oddsidemargin -1cm \evensidemargin -1cm \textwidth 17.5cm
%topmargin 1.2cm \textheight 25cm
%\setlength{\textheight}{25cm}

\def \ds{\displaystyle}
\theoremstyle{definition}
\newtheorem{quest}{Ejercicio}
\thispagestyle{empty}

\begin{document}

\hfill \hfill
\begin{tabular}[c]{|c|c|c|c|c|}
\hline 
1 & 2 & 3 & 4 & Calificaci\'on \\ \hline \hline
\hspace{1 cm} & \hspace{1 cm} & \hspace{1 cm} & \hspace{1 cm} & \hspace{1.3 cm} \\
\hspace{1 cm} & \hspace{1  cm} & \hspace{1  cm} & \hspace{1  cm} & \hspace{1.3 cm} \\ \hline

\end{tabular}

\vspace{-1cm}

\noindent Nombre y apellido:

\vspace{0.2cm}

\noindent Número de libreta:

\noindent \hrulefill 

\smallskip
\renewcommand{\baselinestretch}{1.1}

\begin{center}
	\large \textbf{\nombremateria} 
	
	\numeroparcial -- \fecha
\end{center}

\noindent\textit{Para aprobar el parcial deberá obtenerse un puntaje mayor o igual que 4. \textbf{Entregue todos los ejercicios en hojas separadas.}}
% \vspace{.1cm}
\renewcommand{\theenumi}{\textbf{\arabic{enumi}}}

%%% EMPIEZA EL EJERCICIO 1 %%%
\fontsize{11pt}{0.4cm}\selectfont
\begin{quest} \verb|[2. pts.]| Utilizar la Fase I del Método de las Dos Fases para hallar una solución básica factible inicial del siguiente Problema Lineal:
\[
\begin{array}{rrcl}
\min & \multicolumn{3}{l}{z = 2x_1-2x_2-x_3+x_4} \\
s.a: & x_1+x_2+x_3 & = & 8 \\
 & x_2-x_3+x_4 & \leq & -4 \\
 & -x_2-2x_4 & \geq & -2 \\
 & x_4 & \leq & 0 \\
 & x_1, x_2, x_3 & \geq & 0
\end{array}
\]

\end{quest}

%%% TERMINA EL EJERCICIO 1 %%%

%%% EMPIEZA EL EJERCICIO 2 %%%

\begin{quest} \verb|[1 pt.]|Los siguientes diccionarios corresponden a una iteración de SIMPLEX de diferentes problemas lineales con objetivo de maximizar.

\begin{multicols}{2}
    \begin{itemize}
        \item[1)]
            \[
            \begin{array}{rrrrrrr}
                x_2&=&2&-&\frac{1}{2}x_1&-&x_5\\
                x_3&=&3&-&x_1&+&x_5\\
                x_4&=&\frac{1}{3}&+&x_1&&\\ \hline
                z&=&-3&&&-&\frac{1}{2}x_5
            \end{array}
            \]
        \item[3)]
            \[
            \begin{array}{rrrrrrr}
                x_1&=&1&&&+&\frac{1}{2}x_5\\
                x_4&=&1&&&-&x_3\\
                x_3&=&2&&&-&\frac{4}{3}x_5\\ \hline
                z&=&&-&\frac{1}{4}x_2&-&3x_3
            \end{array}
            \]
        \item[2)]
            \[
            \begin{array}{rrrrrrr}
                x_3&=&7&&&-&x_4\\
                x_1&=&\frac{2}{3}&+&\frac{1}{2}x_2&+&\frac{1}{4}x_5\\
                x_4&=&&&2x_2&-&4x_5\\ \hline
                z&=&-4&+&\frac{3}{4}x_2&-&2x_5
            \end{array}
            \]
        \item[4)]            \[
            \begin{array}{rrrrrrr}
                x_4&=&&&2x_1&+&x_5\\
                x_2&=&1&+&x_1&-&3x_5\\
                x_3&=&\frac{3}{2}&-&x_1&-&x_5\\ \hline
                z&=&2&-&x_1&-&x_5
            \end{array}
            \]

    \end{itemize}
\end{multicols}
Indicar y justificar a qué situación corresponde cada uno de ellos:
\begin{multicols}{2}
\begin{enumerate}[label=\alph*)]
    \item Solución única no degenerada.
    \item Solución única degenerada.
    \item No acotado.
    \item Infinitas soluciones.
\end{enumerate}    
\end{multicols}

\end{quest}
%%% TERMINA EL EJERCICIO 2 %%%

%%% EMPIEZA EL EJERCICIO 3 %%%


\begin{quest} \verb|[4.5 pts.]| 
    Considerar el siguiente Problema Lineal:
    \[
    \begin{array}{rrcl}
    \max & \multicolumn{3}{l}{(3+\beta)x_1+x_2-4x_3+(2+\beta)x_4} \\
    s.a: & 2x_1+x_2-2x_3+x_4 & \geq & 10+\alpha \\
     & -x_2+3x_3-2x_4 & \leq & -8-\gamma \\
     & x_1, x_2, x_3, x_4 & \geq & 0\\
    \end{array}
    \]
    \begin{enumerate}[label=\alph*)]
        \item\verb|[1.5]| Para $\alpha =\beta=\gamma=0$, hallar la solución óptima mediante la solución de su problema dual.
        \item\verb|[1.0]| Para el problema primal, fijando $\alpha=\gamma=0$, hallar todos los valores de $\beta$ para las cuales la base óptima sigue siendo factible.
        \item\verb|[1.5]| Para el problema primal, fijando $\beta=0$, hallar todos los $(\alpha, \gamma)\in\mathbb{R}^2$ para los cuales la solución hallada en el ítem 1 sigue siendo óptima. 
    \end{enumerate}  
\end{quest}

%%% TERMINA EL EJERCICIO 3 %%%

%%% EMPIEZA EL EJERCICIO 4 %%%


\begin{quest} \verb|[2.5 pts.]| Utilizar el algoritmo de Branch \& Bound para resolver el siguiente problema de Programación Lineal Entera, detallando cada división en subproblemas mediante el esquema de árbol visto en clase. 
\[
\begin{array}{rrcl}
\max & \multicolumn{3}{l}{\frac{1}{3}x_1+x_2} \\
s.a: & 2x_1-2x_2 & \geq & -5 \\
 & 2x_1+4x_2 & \leq & 31 \\
 & 2x_1+2x_2 & \leq & 27 \\
 & x_2 & \geq & 0 \\
 & x_1, x_2& \in & \mathbb{Z}\\
\end{array}
\]

\end{quest}
%%% TERMINA EL EJERCICIO 4 %%%

%%% TERMINA EL PARCIAL %%%



\end{document}
